
Supposons maintenant \(X\) noethérien et \(Y\) régulier, ce qui entraine que \(v\) est surjectif (3.7). On a vu dans la démonstration de 3.7 que pour tout \(x'\in K^{\bullet}(X')\) il existe \(y'\in K^{\bullet}(Y')\) tel que :
\[
  x'=f^*(f_*(x'))+j_*(y').
\]
La même propriété reste évidemment vraie après tensorisation par \(\mathbb Q\). Par suite, si \(x'\in\Fil^kK^{\bullet}(X')_{\mathbb Q}\), il existe \((y',x)\in K^{\bullet}(Y')_{\mathbb Q}\times K^{\bullet}(X)_{\mathbb Q}\) tels que :
\[
  x'=v(y',x),\qquad x\in\Fil^kK^{\bullet}(X)_{\mathbb Q}.
\]
En effet, il suffit de prendre \(x=f_*(x')\).

On filtre l'anneau \(K^{\bullet}(Y')_{\mathbb Q}\times K^{\bullet}(X)_{\mathbb Q}\) défini en 4.7 en posant :
\[
  \Fil^k\bigl(K^{\bullet}(Y')_{\mathbb Q}\times K^{\bullet}(X)_{\mathbb Q}\bigr)
  =
  \Fil^{k-1}K^{\bullet}(Y')_{\mathbb Q}\times \Fil^kK^{\bullet}(X)_{\mathbb Q}.
\]
Il résulte immédiatement de (4.7.1) que c'est une filtration d'anneau. Pour montrer que pour tout \(k\) l'homomorphisme :
\[
  \Fil^k\bigl(K^{\bullet}(Y')_{\mathbb Q}\times K^{\bullet}(X)_{\mathbb Q}\bigr)
  \xrightarrow{v}
  \Fil^kK^{\bullet}(X')_{\mathbb Q}
\]
est surjectif, il suffit donc de montrer que si \(x'= \gamma^k(x'_1)\), avec \(\epsilon(x'_1)=0\), il existe \((y',x)\in\Fil^k(K^{\bullet}(Y')_{\mathbb Q}\times K^{\bullet}(X)_{\mathbb Q})\) tel que \(x'=v(y',x)\).

Or d'après ce qui précède, il existe \((y'_1,x_1)\in K^{\bullet}(Y')_{\mathbb Q}\times K^{\bullet}(X)_{\mathbb Q}\) tel que \(x'_1=v(y'_1,x_1)\) et \(x_1\in\Fil^1K^{\bullet}(X)_{\mathbb Q}\). D'après 4.7, on a alors :
\[
  x'=v\bigl(\gamma^k_*(y'_1,x_1)\bigr).
\]
Or :
\[
  \gamma^k_*(y'_1,x_1)
  =
  \left(\sum_{p=1}^{k}\gamma^{k-p}(x_1)\gamma^p(L,y'_1),\;
  \gamma^k(x_1)\right)
\]
(conséquence immédiate de (4.7.2)). Comme \(x_1\) est d'augmentation nulle, \(\gamma^k(x_1)\in\Fil^kK^{\bullet}(X)_{\mathbb Q}\), et compte tenu de V 5.7,
\[
  \sum_{p=1}^{k}\gamma^{k-p}(x_1)\gamma^p(L,y'_1)
  \in\Fil^{k-1}K^{\bullet}(Y')_{\mathbb Q},
\]
d'où le résultat.

